\section{Rational boxes and the sharp inequality in every dimension}
\label{sec:euclidean}

The passage from Brunn--Minkowski to isoperimetry is familiar in classical
geometry \cite{Gardner2002}. Here the central theorem is instead formulated
entirely with rational volumes and an extra rational parameter. This removes
all roots from the finite statement and allows the geometric cuts to be
implemented by bounded rational arithmetic.

\subsection{Finite volume and the box inequality}

A rational box in dimension \(n\geq1\) is a product
\(\prod_{i=1}^n[a_i,b_i]\) with rational \(a_i<b_i\). A box union is
a finite list whose boxes have pairwise disjoint interiors. Its volume
is the sum of their products of side lengths. Boundary conventions do not
affect this volume. To compute the volume of an arbitrary finite union,
sort the endpoints in each coordinate and test the elementary cells of
the resulting grid for inclusion. This is a finite rational procedure.
Inclusion of two unions can be checked on the same common grid.

For two box unions \(X,Y\), their vector sum \(X+Y\) is the union of
all pairwise box sums: add the lower endpoints and the upper endpoints
coordinate by coordinate. These boxes may overlap, so use the preceding
grid procedure to compute \(V(X+Y)\).

\begin{lemma}[A rational product inequality]\label{lem:box-product}
If \(t_1,\ldots,t_n,q\) are nonnegative rationals and
\(\prod_i t_i\geq q^n\), then
\[
 \prod_{i=1}^n(1+t_i)\geq(1+q)^n.
\]
\end{lemma}
\begin{proof}
For \(1\leq j\leq n\), let \(C=\binom nj\) and let \(e_j\)
be the sum of the \(C\) products \(\prod_{i\in J}t_i\), where
\(|J|=j\). Each \(t_i\) occurs in exactly
\(\binom{n-1}{j-1}\) of these products. Thus their product is
\[
 \left(\prod_i t_i\right)^{\binom{n-1}{j-1}}
 \geq q^{n\binom{n-1}{j-1}}=q^{jC}.
\]
Finite arithmetic--geometric mean gives
\((e_j/C)^C\geq q^{jC}\). Since nonnegative rational powers are
order-reflecting, \(e_j\geq Cq^j\). With \(e_0=1\), expand the
left side and sum these inequalities; the binomial theorem gives the
right side. Every exponent, list and subset in this proof is finite.
\end{proof}

If \(X,Y\) are single boxes with side lengths \(a_i,b_i>0\), put
\(t_i=b_i/a_i\). Whenever \(q^n V(X)\leq V(Y)\), the lemma gives
\begin{equation}\label{eq:box-base}
 V(X+Y)\geq V(X)(1+q)^n.
\end{equation}
This is the only numerical inequality needed for the general finite case.

\subsection{Matching rational fractions by explicit cuts}

\begin{lemma}[Rational volume cut]\label{lem:rational-cut}
Let \(Y\) be a nonempty box union, fix a coordinate \(i\), and let
\(0<\lambda<1\) be rational. A rational \(t\) can be computed such
that the pieces \(Y^- =Y\cap\{x_i\leq t\}\) and
\(Y^+=Y\cap\{x_i\geq t\}\) have volumes
\(\lambda V(Y)\) and \((1-\lambda)V(Y)\).
Each list contains at most as many boxes as the original list.
\end{lemma}
\begin{proof}
Sort the finitely many endpoints in coordinate \(i\).
Between two consecutive endpoints \(u<v\), the cross-sectional
\((n-1)\)-volume is a constant rational \(A\geq0\); for \(n=1\)
the empty product used here is one. Scan these intervals, adding
\(A(v-u)\), until the target \(\lambda V(Y)\) is reached or crossed.
If it is reached at an endpoint, use that endpoint. Otherwise the target
lies strictly inside an interval with \(A>0\); if the volume accumulated
at its left endpoint is \(F\), choose
\(t=u+(\lambda V(Y)-F)/A\).
All operations are rational and the scan length is bounded by the number
of endpoints. Clipping each original box produces at most one nonempty
box in either specified half-space. It does not subdivide a box into
two pieces \emph{within the same child list}.
\end{proof}

\begin{theorem}[Brunn--Minkowski with a rational parameter]
\label{thm:rational-bm}
Let \(X,Y\) be nonempty finite unions of rational boxes with disjoint
interiors, in dimension \(n\geq1\), and let \(q\geq0\) be rational.
Then
\begin{equation}\label{eq:rational-bm}
 q^n V(X)\leq V(Y)
 \quad\Longrightarrow\quad
 V(X+Y)\geq V(X)(1+q)^n.
\end{equation}
\end{theorem}
\begin{proof}
We describe a recursion on \(N\), the sum of the two list lengths.
When both lists have length one, use \eqref{eq:box-base}.
Suppose first that \(X\) has at least two boxes. Because their interiors
are disjoint, two selected boxes have disjoint open coordinate intervals
in some coordinate \(i\). Choose a rational hyperplane \(x_i=s\)
between them, so one whole selected box is on each side. Clip every box
of \(X\) at this hyperplane, discard zero-width pieces, and call the
results \(X^-,X^+\). Both have positive volume. Put
\(\lambda=V(X^-)/V(X)\), so \(0<\lambda<1\).
Use Lemma~\ref{lem:rational-cut} to cut \(Y\) at \(x_i=t\) with
the same volume fraction. The identities are
\[
 V(X^-)=\lambda V(X),\quad V(Y^-)=\lambda V(Y),\qquad
 V(X^+)=(1-\lambda)V(X),\quad V(Y^+)=(1-\lambda)V(Y).
\]
Hence the same hypothesis \(q^nV(X^\pm)\leq V(Y^\pm)\) holds
in both children.

If \(X\) originally had \(m\) boxes and \(Y\) had \(r\), each
\(X\) child has at most \(m-1\) boxes: the selected whole box on
the other side is absent, and each remaining original contributes at
most one piece. Each \(Y\) child has at most \(r\) boxes.
Thus the sum of list lengths in \emph{each recursive call} is at most
\(N-1\). The sum of the sizes of both calls need not decrease, and is
not used as the induction measure.

The induction hypotheses give
\[
 V(X^-+Y^-)\geq V(X^-)(1+q)^n,
 \qquad V(X^++Y^+)\geq V(X^+)(1+q)^n.
\]
The first sum is in \(\{x_i\leq s+t\}\), the second in
\(\{x_i\geq s+t\}\). Their interiors are disjoint and both lie in
\(X+Y\), so finite additivity of box volume allows their lower bounds
to be added. This proves \eqref{eq:rational-bm}.

If \(X\) has one box but \(Y\) has several, choose the separating cut
in \(Y\) and match its volume fraction in \(X\). The same equations
and the same decrease of \(N\) apply. This exhausts the cases without
ever changing the direction of the hypothesis involving \(q\).
\end{proof}

\paragraph{Explicit arithmetic bound.}
The recursion has depth at most \(N-2\) and fewer than \(2^{N+1}\)
nodes. It must retain the clipped lists, rather than replacing each child
by a fresh grid decomposition. Suppose endpoint numerators and positive
denominators initially have absolute value at most \(2^B\).
Taking an unnecessarily large bound, set
\[
 B_0=B+2,\qquad
 B_{j+1}=2^{\,100(n+1)(N+1)^2(B_j+1)^2}.
\]
After \(j\) levels, all endpoint numerators and denominators are bounded
by \(2^{B_j}\). To verify this, note that a cut requires only products
of at most \(n\) endpoint differences, sums of at most \(N\) such
products, and a fixed number of rational divisions; forming a common
denominator multiplies at most this many previous denominators. The
displayed bound dominates these elementary bit-length bounds, even
without cancellation. Evaluation of the sum volume at the root uses
at most \((2mr)^n\) elementary grid cells. These are primitive recursive
bounds uniform in all the inputs. A postorder traversal with these bounds
gives a bounded Boolean induction invariant, as specified in
Section~\ref{sec:foundations}.

\subsection{Recovering the conventional sharp statement}

Let \(a=V(X)>0\), \(b=V(Y)>0\). By prescribed bisection choose
\(q_j\geq0\) rational with
\(q_j^na\leq b\) and \(0\leq(b/a)^{1/n}-q_j\leq2^{-j}\).
Theorem~\ref{thm:rational-bm} holds for every \(q_j\).
The factorization bound for powers supplies an explicit error budget
for \(a(1+q_j)^n\); therefore the order definition gives
\[
 V(X+Y)\geq(a^{1/n}+b^{1/n})^n,
 \qquad V(X+Y)^{1/n}\geq a^{1/n}+b^{1/n}.
\]
Replacing \(X,Y\) by positive rational dilates proves the weighted
version. Named positive dilation factors follow by their rational error
bounds. This passage uses no convergence theorem without a modulus.

\subsection{A primitive recursive volume name for the unit ball}

We need a named constant \(\omega_n\), not a postulated Lebesgue measure
of an unencoded set. Take the grid of side \(h=2^{-m}\) in \([-2,2]^n\).
Let \(B_m^-\) be the union of cells whose farthest corner from the origin
has squared norm at most one. Let \(B_m^+\) be the union of cells whose
minimum squared norm is at most one. In the latter test, the minimizing
coordinate is the endpoint nearest zero, or zero if the coordinate
interval contains zero. Both tests are rational. The unions enclose
the unit ball from inside and outside, respectively, and refine
monotonically under dyadic subdivision, up to their shared boundaries.

Here is an explicit estimate independent of integration. Suppose
\(nh\leq1/4\), and put
\(c=(1-2nh)/(1+nh)>0\).
Every point of \(B_m^+\) has norm at most \(1+nh\), since a cell has
diameter at most \(nh\). Thus every point of \(cB_m^+\) has norm at
most \(1-2nh\). The grid cell containing such a point has every corner
of norm at most \(1-nh\), so it belongs to \(B_m^-\). Hence
\[
 cB_m^+\subseteq B_m^-\subseteq B_m^+,
 \qquad
 0\leq V(B_m^+)-V(B_m^-)
 \leq (1-c^n)4^n\leq3n^2 4^n h.
\]
The last inequality uses \(1-c^n\leq n(1-c)\) and
\(1-c\leq3nh\). All volume comparisons are comparisons of finite
rational box unions; the norm estimates use the finite Cauchy--Schwarz
inequality.

Define \(\omega_n\) by these nested intervals. For example
\(m=k+3n^24^n+4n+4\) is more than enough for interval width at most
\(2^{-k}\). The contained cube \([-1/n,1/n]^n\) and the bounding cube
give positive numerical bounds
\[
 (2/n)^n\leq\omega_n\leq 2^n.
\]
The latter upper bound follows directly from
\(B_m^+\subseteq[-1-h,1+h]^n\), hence
\(V(B_m^+)\leq(2+2h)^n\); the binomial theorem gives the explicit
vanishing error as \(h\) decreases. Rational dilation of the
finite approximants proves that a ball of radius \(R>0\) has volume
\(\omega_nR^n\). In dimension two the disk area calculation in the
planar appendix identifies \(\omega_2=\pi\).

\subsection{Precisely stated regularity data for outer volume}

An \emph{effective outer-volume presentation} of a region consists of
the following data, all supplied by fixed primitive recursive families.
\begin{enumerate}
\item Named nonnegative numbers \(V,P\), an integer bound \(M\geq1\)
with \(V,P\leq M\), and a positive rational \(v_0\leq V\).
\item Inner rational box unions \(X_k\) with
\(0\leq V-V(X_k)\leq2^{-k}\).
\item For rational \(t>0\) and precision \(k\), rational box unions
\(Z_{t,k}\) which contain
\(X_k+tB_k^-\), and a named outer volume \(W(t)\) with
\(V(Z_{t,k})\leq W(t)+2^{-k}\).
Here and below \(B_k^-\) abbreviates the unit-ball approximant at the
mesh index \(m(k)=k+3n^24^n+4n+4\), so its volume error is at most
\(2^{-k}\). A larger explicitly prescribed mesh index is also allowed.
\item A primitive recursive modulus \(\nu\) such that
\begin{equation}\label{eq:tube-modulus}
 0<t\leq2^{-\nu(j)}\quad\Longrightarrow\quad
 W(t)\leq V+t(P+2^{-j}).
\end{equation}
\end{enumerate}
The containments can be supplied as finite common-grid certificates.
If an underlying geometric set \(S\) is also specified, the intended
meaning is \(X_k\subseteq S\),
\(W(t)=V(S+tB_1)\), and \(P\) is its boundary area or its outer
Minkowski perimeter. The finite-data theorem only uses the displayed
volume and containment requirements. In particular it never assumes an
isoperimetric inequality in the input.

Condition \eqref{eq:tube-modulus} supplies the rate of the one-sided
derivative calculation. It is extra constructive regularity data. The
report does not assert that an arbitrary classical \(C^1\) boundary,
given without any effective presentation, automatically supplies it.

\begin{theorem}[Sharp Euclidean isoperimetry with outer-volume data]
\label{thm:euclidean-isoperimetry}
For every presentation just specified, in every \(n\geq2\),
\[
 P\geq n\omega_n^{1/n}V^{(n-1)/n},
 \qquad
 P^n\geq n^n\omega_n V^{n-1}.
\]
\end{theorem}
\begin{proof}
First fix rational \(t>0\). Apply the finite box theorem to
\(X_k,tB_k^-\) and use their inclusion in \(Z_{t,k}\).
For sufficiently large prescribed \(k\), \(V(X_k)\geq v_0/2>0\).
The preceding root estimates give, at every requested precision,
\[
 W(t)\geq\bigl(V^{1/n}+t\omega_n^{1/n}\bigr)^n.
\]
For clarity, this is a finite error calculation: if both input volume
errors are at most \(\eta\), their root errors are at most
\(\eta^{1/n}\); if all roots and their sums are bounded by \(D\),
the error in their \(n\)-th powers is at most
\(nD^{n-1}(1+t)\eta^{1/n}\), plus the outer-volume error \(\eta\).
Use \(D=2(M+4^n+1)(1+t)\), for example, and choose a dyadic
\(\eta\) making this expression smaller than the desired tolerance.
All choices are primitive recursive in the numerical data and precision.

Discarding the nonnegative terms of degree at least two in the binomial
expansion now gives
\[
 W(t)\geq V+tn\omega_n^{1/n}V^{(n-1)/n}.
\]
For an arbitrary \(j\), choose the specific rational
\(t=2^{-\nu(j)-1}>0\) and compare with
\eqref{eq:tube-modulus}. Subtract \(V\), divide by this known positive
rational, and obtain
\[
 n\omega_n^{1/n}V^{(n-1)/n}\leq P+2^{-j}.
\]
Because this holds at every precision, it is the asserted order
inequality. Raising to the positive integer power \(n\) gives the
second form. This proves the theorem without an unbounded selection of
small radii or an unquantified limiting argument.
\end{proof}

There is also a particularly useful finite numerical trace of this proof.
Suppose the box data have volumes \(a,b,c\), \(q^na\leq b\),
\(X+tB\subseteq Z\), and
\[
 c\leq v+t(p+\varepsilon),\qquad 0\leq v-a\leq t\varepsilon.
\]
Then the rational box theorem gives
\(c\geq a(1+tq)^n\geq a+ntaq\), and therefore
\begin{equation}\label{eq:finite-isop-trace}
 p+2\varepsilon\geq naq.
\end{equation}
This is a convenient certificate target for a proof checker.

\subsection{Examples verifying the regularity contract}

\paragraph{Balls and sharpness.}
For a ball of radius \(R>0\), the parallel set is the ball of radius
\(R+t\). Therefore
\[
 V=\omega_nR^n,\qquad P=n\omega_nR^{n-1},\qquad
 W(t)=\omega_n(R+t)^n.
\]
For \(0<t\leq1\), the remaining terms of the binomial expansion are
at most
\[
 t^2\,\omega_n\sum_{j=2}^n\binom nj(R+1)^{n-j}.
\]
A rational upper bound on this coefficient gives a modulus \(\nu\)
by direct dyadic comparison. Inner and outer grid approximants were
constructed above. Substitution gives
\(P^n=n^n\omega_nV^{n-1}\), proving the sharpness of the universal
constant. This establishes that balls attain equality; it does not claim
a complete higher-dimensional classification of every equality set.

\paragraph{A rational rectangular box.}
For side lengths \(a_1,\ldots,a_n>0\), take
\[
 V=\prod_i a_i,\qquad P=2\sum_i\prod_{j\ne i}a_j.
\]
The Euclidean \(t\)-parallel set is contained in the rectangular box
with sides \(a_i+2t\). Thus it is enough to take
\(W(t)=\prod_i(a_i+2t)\), even though this is only an upper envelope
for the actual Euclidean parallel volume. For \(0<t\leq1\),
\[
 W(t)\leq V+tP+t^2K,\qquad
 K=\sum_{\substack{J\subseteq[n]\\|J|\geq2}}
 2^{|J|}\prod_{i\notin J}a_i.
\]
Choose \(\nu(j)\) so that \(K2^{-\nu(j)}\leq2^{-j}\).
The inner box can be the box itself and the required outer boxes are
explicit. This verifies the full contract without differential geometry.

\paragraph{Every finite rational box union.}
The same conclusion holds for a nonempty finite union \(X\) of rational
boxes. Here \(P\) is the total area of its exposed coordinate faces,
computed after subdividing into a common endpoint grid. Each face between
an occupied and an unoccupied grid cell is counted once; shared faces of
two occupied cells are not counted. This is a nonnegative rational number.

For a direct finite proof of the tube bound, expand each box by
\([-t,t]^n\) and take the union. This contains the Euclidean
\(t\)-parallel set. If the original boxes have lower and upper endpoints
\(a_{ji},b_{ji}\), finite inclusion--exclusion gives its exact volume
\[
 W(t)=\sum_{\varnothing\ne J\subseteq\{1,\ldots,m\}}
 (-1)^{|J|+1}
 \prod_{i=1}^n
 \max\!\left(0,\min_{j\in J}b_{ji}-\max_{j\in J}a_{ji}+2t\right).
\]
Let \(\delta>0\) be the least nonzero distance between two distinct
endpoints in the same coordinate, found by a finite rational scan, and
put \(t_0=\min(1,\delta/4)\). Such a gap exists because every box has
positive side lengths. For \(0<t\leq t_0\), every strictly negative
gap in the displayed formula stays negative, while each zero gap
contributes \(2t\). Thus this formula is one explicitly computable
polynomial in \(t\) throughout that interval.

Its constant coefficient is \(V(X)\), and its linear coefficient is
exactly \(P\). To verify the latter assertion without a differential
argument, first replace the boxes by their occupied common-grid cells;
Minkowski addition distributes over the finite union, so the expanded
union is unchanged. The expansion of one cell contributes its total
face area to the linear coefficient. Two cells sharing a face of area
\(F\) have an expanded intersection of volume \(2tF+O(t^2)\), so
their inclusion--exclusion term cancels the two copies of that internal
face. Any other nonempty intersection involves at least two independent
zero coordinate gaps, and therefore has no constant or linear term.
Intersections across a positive grid gap are empty for \(t\leq t_0\).
These assertions follow by multiplying the interval-length factors in
the same finite formula. Only the exposed faces remain.

Write the resulting polynomial as
\(W(t)=V(X)+tP+\sum_{j=2}^n c_jt^j\), and compute the rational
\(K=\sum_{j=2}^n|c_j|\). Then
\[
 W(t)\leq V(X)+tP+Kt^2\qquad(0<t\leq t_0).
\]
For example, take an integer \(C\geq1+K+t_0^{-1}\) and set
\(\nu(j)=j+C+2\). Since \(2^C\geq C\), this guarantees both
\(2^{-\nu(j)}\leq t_0\) and
\(K2^{-\nu(j)}\leq2^{-j}\).
Use \(X_k=X\) for every \(k\), and take \(Z_{t,k}\) to be the
explicit expanded box union, subdivided into a disjoint grid if needed.
The four regularity conditions follow. Consequently the Euclidean
isoperimetric inequality in Theorem~\ref{thm:euclidean-isoperimetry}
is unconditional for these finite rational box unions with their
exposed-face perimeter.

\paragraph{The zero-volume case.}
If a presentation comes with a proof \(V=0\), the polynomial inequality
for \(n\geq2\) follows from \(P\geq0\). We did not decide whether
an arbitrary named real volume is zero. The positive-volume theorem uses
the supplied lower bound, as a constructive presentation of a nonempty
open region normally does.

\paragraph{What is and is not covered.}
The theorem applies to effective smooth or polyhedral domains whenever
the four listed approximation and tube requirements have been proved
for their presentation. It is a fully proved implication from explicit
geometric regularity data. A general construction of those data from
arbitrary \(C^1\) charts, and equivalence with every classical definition
of perimeter or Hausdorff measure, are separate tasks not concealed in
this implication. The article's assertion for arbitrary sets with a
liminf Minkowski boundary is not claimed in that unrestricted form here.

